Proxy variables and the challenges of using them correctly have long been considered in the causal inference literature \citep{wickens1972note,frost1979proxy}. Understanding what is the best way to derive and measure possible proxy variables is an important part of many observational studies \citep{filmer2001estimating,kolenikov2009socioeconomic,wooldridge2009estimating}.
Recent work by \citet{cai2008identifying,greenland2008bias}, building on the work of \citet{greenland1983correcting,selen1986adjusting}, has studied conditions for causal identifiability using proxy variables.
The general idea is that in many cases one should first attempt to infer the joint distribution $p(\*X,\*Z)$ between the proxy and the hidden confounders, and then use that knowledge to adjust for the hidden confounders \citep{wooldridge2009estimating,pearl2012measurement,kuroki2014measurement,miao2016identifying,edwards2015all}.
For the example in Figure~\ref{fig:hid1}, \citet{cai2008identifying,greenland2008bias,pearl2012measurement} show that if $\*Z$ and $\*X$ are categorical, with $\*X$ having at least as many categories as $\*Z$, and with the matrix $p(\*X,\*Z)$ being full-rank, one could identify the causal effect of $\*t$ on $\*y$ using a simple matrix inversion formula, an approach called ``effect restoration''.
Conditions under which one could identify more general and complicated proxy models were recently given by \citep{miao2016identifying}. 



